\documentclass[10pt,a4paper]{amsart}
\usepackage[margin=19mm]{geometry}
\usepackage{amsmath,amssymb,mathtools}
\usepackage[scr=rsfs,cal=euler]{mathalfa}
\usepackage{xfrac}
\usepackage[unicode,hidelinks,bookmarks=false]{hyperref}
\newcommand{\ie}{i.e.\ }
\newcommand{\cf}{cf.\ }
\newcommand{\resp}{respectively\ }
\newcommand{\Subsection}{\subsection}
\providecommand{\nobreakdash}{\nobreak\hbox{-}}
\setlength{\emergencystretch}{2em}
\multlinegap0pt
\usepackage{amssymb}
\usepackage{float}
\newtheorem{lemma}{Lemma}
\title{Covering by Centralizers}
\author{Mark L. Lewis, Ryan McCulloch}
\date{Source: arXiv:2512.02211v2 (CC BY 4.0)}
\begin{document}
\maketitle

\noindent\emph{Source and licence.} Covering by Centralizers. Version arXiv:2512.02211v2 (CC BY 4.0), distributed by arXiv under the Creative Commons Attribution 4.0 International licence, http://creativecommons.org/licenses/by/4.0/, as recorded in the arXiv licence metadata for this exact version and shown on the abstract page https://arxiv.org/abs/2512.02211v2. Full licence notice: \texttt{licenses/arxiv-2512.02211-CC-BY-4.0.txt}.\par\smallskip

\noindent\textit{Editorial scope.} Counted unit: the article's lemma 2 (source0.tex lines 253--261). The article's complete original proof of it is delivered with it (source0.tex lines 263--273, one \texttt{proof} environment of 11 lines). No mathematics is written, repaired or re-derived: the statement and the proof are the article's own lines, in the article's own notation, and no retained line is substituted or re-flowed.  Retained uncounted as the source-local context the proof needs: lines 182--188; lines 201--203.  Nothing was omitted from any retained span, so the package carries no omission record.  No retained line contains a citation, so the wrapper carries no bibliography and no bibliography entry.  No retained line contains a figure environment, an image-inclusion command or drawing code, so no image is delivered and the package declares no asset dependency.  Editorial formatting only, in addition: an AMS \texttt{amsart} preamble that restates the article's own declarations and macros used by the retained lines, namely the environments \texttt{lemma} on separate counters, together with the standard packages those lines need.  The retained environments print in this document as Lemma 1 in order of appearance, whereas the article prints the same items with the numbering of its own full document.  The complete native source of the article as distributed in the arXiv e-print for this exact version is bundled unchanged as \texttt{sources/190183/source0.tex}.
\begin{lemma} \label{star 1}
Let $G$ be a group and let $g, h \in G \setminus Z(G)$.  Then 
\begin{enumerate}
\item $C_G (g) \le C_G (h)$ if and only if $Z(h) \le Z(g)$.
\item $C_G (g) = C_G (h)$ if and only if $Z (g) = Z(h)$.
\end{enumerate}
\end{lemma}

\begin{lemma} \label{intersection}
Let $G$ be a group and let $g, h \in G \setminus Z(G)$.   Then either $Z(g) \cap Z(h) = Z(G)$ or $(Z(g) \cap Z(h)) = \prod Z(c)$ where $c$ runs over the elements in $(Z (g) \cap Z(h)) \setminus Z(G)$.
\end{lemma}

\begin{lemma} \label{star 2}
Let $G$ be a nonabelian group and $g \in G \setminus Z(G)$.  If $g_1, \dots, g_t \in G \setminus Z(G)$ such that $C_G (g_1), \dots, C_G (g_t)$ are distinct and comprise all of the elements of ${\mathcal C} (G)$ that properly contain $C_G (g)$, then $\displaystyle Z(g) \setminus Z(G) = Z^* (g) \cup (\bigcup_{i=1}^t Z^* (g_i))$.   In particular, $\displaystyle \bigcup_{i=1}^t Z^* (g_i)$ is a proper subset of $Z(g) \setminus Z(G)$. In addition, if $h \in G \setminus Z(G)$, then the following are equivalent:
\begin{enumerate}
\item $h \in Z(g)$.
\item $Z^* (h) \subseteq Z(g)$.
\item $Z(h) \le Z(g)$.
\item $C_G (g) \le C_G (h)$.
\end{enumerate}
\end{lemma}

\begin{proof}
Since $C_G (g) < C_G (g_i)$, we have $Z^* (g_i) \subseteq Z(g_i) \setminus Z(G) \subseteq Z(g) \setminus Z(G)$ by Lemma \ref{star 1} for each $i$.  It follows that $\displaystyle \bigcup_{i=1}^t Z^* (g_i) \subseteq Z(g) \setminus Z(G))$.  Since $Z^* (g) \subseteq Z(g) \setminus Z (G)$, we have $\displaystyle Z^* (g) \cup (\bigcup_{i=1}^t Z^* (g_i)) \subseteq Z(g) \setminus Z(G)$.
	
On the other hand, if $h \in Z(g) \setminus Z(G))$, then $h$ centralizes $C_G (g)$.  We see that $C_G (g)$ centralizes $h$; so $C_G (g) \le C_G (h)$.  If $C_G (g) = C_G (h)$, then $h \in Z^* (g)$.  Thus, we may assume $C_G (g) < C_G (h)$.  Now, $C_G (h)$ is an element of ${\mathcal C} (G)$ that contains $C_G (g)$.   Since $h \not\in Z(G)$, it follows that $C_G (h) = C_G (g_i)$ for a unique $i$ with $1 \le i \le t$.  This implies $h \in Z^* (g_i)$, and so, $\displaystyle h \in \bigcup_{i=1}^t Z^* (g_i)$.  We conclude that $\displaystyle Z(g) \setminus Z(G)) = Z^* (g) \cup (\bigcup_{i=1}^t Z^* (g_i))$.   Since the $Z^*$'s are disjoint and $Z^* (g)$ is not empty, this implies that  the conclusion that $\displaystyle \bigcup_{i=1}^t Z^* (g_i)$ is proper in $Z(g) \setminus Z(G)$ is now immediate. 
	
We know that $Z(h) \le Z(g)$ is equivalent to $C_G (h) \le C_G (g)$.  Suppose $C_G (g) \le C_G (h)$.  If $C_G (g) = C_G (h)$, then $Z^* (g) = Z^* (h)$.  Also, when $C_G (g) < C_G (h)$, we have shown that $Z^* (h) \subseteq Z (g)$.  This shows that $Z^* (h) \subseteq Z^* (g)$ in both cases.
	
It suffices to show that if $Z^* (h) \subseteq Z (g)$, then $Z(h) \le Z(g)$.  We show the contrapositive.  Suppose $Z (h) \not\le Z (g)$.  This implies that $h \not\in Z(g)$ since by Lemma \ref{intersection}, we have $Z(h) \cap Z(g)$ is the product of $Z (c)$ where $c$ runs over the elements in $Z(h) \cap Z(g)$.  Since $h$ lies in $Z^* (h)$, this implies that $Z^* (h)$ is not contained in $Z(g)$, and we have proved the contrapositive.
	
Suppose $h \in Z(g)$.  Then $\displaystyle h \in Z^* (g) \cup (\bigcup_{i=1}^t Z^* (g_i))$.  Either $h \in Z^* (g)$ or $h \in Z^* (g_i)$ for some $i$.  In either case, $Z (h) \le Z(g)$.  Conversely, if $Z(h) \le Z(g)$, then $h \in Z(h)$ implies $h \in Z(g)$.  This shows that all of the equivalences hold.
\end{proof}

\end{document}
